\section{Challenge guidance and downloadable practice inputs}
\label{learning-guidance}

The November 2026 usability repair addresses a generator defect: broad curriculum
goals had been presented as executable lab instructions, often with the same
course link and no identified dataset. The repair supplies per-track setup,
concrete starter/topic recipes, named downloads, specified changed cases and a
way to save the learner's own observations. Short overview summaries remain
separate from the detailed challenge workspace.

\subsection{Authoring, generation and delivery}
\path{docs/curriculum/specialist-paths.json} owns stable IDs, prerequisites,
concept scenarios and assessment objectives. The companion
\path{docs/curriculum/lab-guides.json} contains authored setup and recipes for
nineteen specialist tracks. Each has a starter, six topic recipes and a specified
independent variation. These files are static authoring data; the application
does not ask a language model to invent a lab when the learner opens it.

\path{scripts/build_specialist_tracks.py} merges these sources into runtime TOML,
matching public exports and catalog/scenario/assessment indexes. Task-specific
guidance overrides may change learner-facing text and duration, but cannot change
IDs, dependencies or verifiers. The Python HTTP walkthrough remains an explicitly
authored override. Existing IDs preserve navigation history, although an old
completion does not retroactively demonstrate a newly clarified skill.

Original synthetic input files live in
\path{tracks/cyber-tracks/materials/}, separated by track. The shared report helper
is included in each ZIP by \path{scripts/package_learning_materials.py}. Archives
use deterministic names, timestamps and permissions and are served from
\path{web/training/}. They contain readable source and training data; downloading
does not execute them. The existing recursive web packaging includes these assets
in distributable installations. The intentionally incomplete notes API listens
only on loopback and uses synthetic identities rather than real authentication.

\subsection{Model and interface contract}
Optional node fields \texttt{preview}, \texttt{requirements}, \texttt{setup}, \texttt{resource\_label} and
\texttt{resource\_note} are validated strings. \texttt{materials} is a list of
objects containing a nonempty label and URL. URLs must use HTTP(S) or the bounded
\path{/training/} asset namespace; arbitrary local paths, script URLs and path
traversal are rejected. Older packs default to empty guidance fields.

The web presentation layer returns previews from overview/list endpoints and full
guidance only from a started or historically completed challenge's workspace.
\texttt{showNode} presents the objective and prerequisite requirements before
Start; the started workspace reveals detailed setup, download links, an ordered
``What to do'' sequence and an explicitly labeled optional reference.
\texttt{instructionBlocks} recognizes fenced code in steps and constructs DOM
text nodes. It does not insert command text as HTML or execute it. The reusable
\texttt{codeBlock} component supplies keyboard-accessible code and a clipboard
button. Long code is constrained within the responsive workspace.

\subsection{Evidence is not an answer key}
The downloaded \path{report.py} prompts for scope, setup, prediction, exact
actions, observations, comparison, explanation, limitations and assistance. It
creates UTF-8 JSON using exclusive file creation. Existing reports are preserved;
empty answers or cancellation produce no report. The learner must execute the
helper and supply their own observations. No generated answer or automatic
completion is included.

The file verifier still checks the server filesystem and optional size/pattern,
not the quality of reasoning. A file saved solely on a learner laptop is invisible
to a hosted/container server until placed at that server's configured lab path.
There is no evidence-upload endpoint in this repair. Artifact/self-reviewed tasks
retain their existing no-automatic-skill-XP behavior. Continuous timers and final
submission rules for native exercises are unchanged.

\subsection{Verification and limits}
\path{scripts/audit_learning_content.py} scans all 28 community packs and 1,002
expanded nodes for known vague placeholders, unnamed materials, misleading
resource roles, repeated introductions and unexplained JSON evidence. It also
checks that bundled download targets exist. A clean diagnostic is not a score for
pedagogy, a learner-outcome study or proof that every external environment works.
CI checks the saved audit, generated packs and deterministic archives.

Regression tests execute the first HTTP recipe and local notes API; inspect
marker cleanup, the permission model, timestamp controls, benign source and mock
assistant; and verify report creation without overwrite. Synthetic PCAPs are
generated without network traffic and decoded with Wireshark when available.
Browser checks cover setup, ZIP downloads, copying, narrow/wide layout and session
reload using disposable profiles.

Windows AD/AD CS, Android, provider IAM, Kubernetes, Security Onion, Volatility and
physical-device environments were not all executed in this repair. Their tasks
state prerequisites and distinguish simulation, preparation and operational
results. In particular, a directory graph is not a Windows effective-access test,
a beacon is not a captured handshake, and a mock assistant is not a real-model
security evaluation. The full authoring checklist and execution limits are in
\path{docs/curriculum/challenge-authoring.md}.

Installed private pack copies require an explicit content update. Before a live
update, back up the database and copies, inspect active attempts and compare
protected learner state afterward. Changed fingerprints reset review/trust and
can invalidate unfinished generated exercises. Saved progress must not be
silently reset as part of a wording repair.

\subsection{Planned disposable terminal boundary}
The comparison of OverTheWire, pwn.college, TryHackMe and CTFd is recorded with
source links and limits in \path{docs/design/isolated-lab-research.md}. OverTheWire's
public SSH model does not establish per-attempt containers. The inspected
pwn.college source uses runc for ordinary challenges and Kata for challenges
marked privileged; that distinction and its deployment-specific mounts matter.
An embedded terminal renderer such as xterm.js is not an isolation boundary.

The selected design separates authenticated attempt state in Ascend from an
allowlisted runner. It requires a gVisor runtime or suitable VM boundary, no host
mounts or secrets, bounded resources, explicit network restrictions, validated
evidence collection and crash-safe teardown. Provisioning must finish before
timing starts; reload must reconnect to the same attempt. Kernel, AD and hardware
tasks require a different capability class from a small local Python lab.

This terminal feature is not implemented or enabled. The inspected workstation
has ordinary Docker runtimes, no gVisor/Kata installation, and no provisioned
runner VM. The preview change does not launch containers or substitute a host
shell. Administrator runtime preparation and deployed isolation/lifecycle tests
are prerequisites for claiming that feature works securely.
